\begin{toreview}

	\section{Severity Scoring Methodology}\label{sec:severity_scoring_methodology}

	Raw article counts from the \gls{slr} conflate two different questions: how much a topic is written about, and how dangerous a topic is. A category can dominate the literature because it is easy to publish on, not because it carries the greatest risk. To separate these axioms, each concept-matrix category is assigned a \emph{severity score}, which is later compared against its \emph{coverage rank} (article count) in Section~\ref{subsec:results_quantitative}.

	\subsection*{Severity Scoring Dimensions}

	The four scoring dimensions are adapted from the ``key risks'' assessment criteria used in IPCC Working Group II reports, first formalized by Smith et al. \cite{smith2009assessing} and later refined by Oppenheimer et al. \cite{oppenheimer2014emergent} and O'Neill et al. \cite{oneill2017ipcc}, who evaluate climate risks along magnitude, timing, persistence/irreversibility, distributional scope, and likelihood/confidence. This work retains the four criteria most transferable to \gls{ai} risk categories --- magnitude, certainty, irreversibility, and scope --- and drops timing and adaptive-capacity, which are climate-specific. This four-dimension structure also echoes the magnitude/scope/irreversibility triad used in existential risk assessment by Bostrom \cite{bostrom2013existential}, giving the framework precedent within both the climate-risk and \gls{ai}-safety literatures.

	\begin{itemize}
		\item \textbf{Magnitude (M)}: ceiling of harm if the failure mode materializes (e.g.\ the Treacherous Turn, \cite{bostrom2014superintelligence}).
		\item \textbf{Certainty (C)}: extent of current empirical evidence, as opposed to purely theoretical risk.
		\item \textbf{Irreversibility (I)}: whether the harm is correctable after the fact.
		\item \textbf{Scope (S)}: breadth of population or systems exposed.
	\end{itemize}

	\subsection*{Rating Process}

	Five independent raters (domain experts) score each category on all four dimensions. The full matrix of raw scores provided by all five raters is available in the Appendix (Table~\ref{tab:rater_scores}). Inter-rater agreement is reported via Krippendorff's $\alpha$ \cite{krippendorff2018content, hayes2007answering}:

	\begin{enumerate}
		\item $\alpha = 1 \to \text{perfect agreement}$
		\item $\alpha = 0 \to \text{agreement no better than chance}$
		\item $\alpha < 0 \to \text{systematic disagreement, worse than random}$
	\end{enumerate}

	\subsection*{Scope Definition for Raters}

	Before scoring, all raters were given the following written scope, to ensure severity judgments reference a common referent rather than each rater's private notion of ``AI'':

	\begin{quote}
	\textbf{Scope of ``AI'' for this severity assessment.} Assess the potential harm of each category by considering both:

	\begin{enumerate}
		\item \textbf{Deployed and near-term systems}: current and near-future AI systems (roughly the next 0--3 years), including \glspl{llm} and other machine learning systems already in use. Base the assessment on documented incidents or failures that are reasonably likely given existing evidence and literature.

		\item \textbf{Frontier and projected systems}: AI systems that could plausibly emerge through continued improvements to current approaches (roughly a 3 to 15 year horizon, including AGI-adjacent capabilities). Base the assessment on technically grounded failure modes discussed in the research literature, such as deceptive alignment, mesa-optimization, and treacherous turns.
	\end{enumerate}

	Do not consider purely speculative or fictional scenarios that lack support in the current technical literature, such as generic science-fiction robot uprisings. When a category applies to both near-term and frontier systems (as is often the case for \emph{Value Alignment}), consider both perspectives when assigning its severity.
	\end{quote}

\end{toreview}